% Môn: Vật lý
% Lớp: 11
% Chương: Chương I. Dao động
% Bài: L11C1 Bài 2. Mô tả dao động điều hòa · Xác định các đại lượng đặc trưng của dao động điều hòa từ phương trình li độ
% ===== Câu 39 =====
% ID: L11C1B2-57-TLN
% Mức: TH
\dangbt{Xác định các đại lượng đặc trưng của dao động điều hòa từ phương trình li độ}
\begin{ex}
% ID: L11C1B2-57-TLN
% Mức: TH
Một vật thực hiện 12 dao động trong $6\,\text{s}$. Tần số dao động theo đơn vị Hz bằng bao nhiêu?
\shortans{2}
\loigiai{
Tần số dao động (f) được tính bằng số dao động thực hiện được trong một đơn vị thời gian.
Công thức: $f = \frac{\text{Số dao động}}{\text{Thời gian}}$
Thay số: $f = \frac{12}{6} = 2\,\text{Hz}$.
}
\end{ex}

% ===== Câu 53 =====
% ID: L11C1B2-58-TLN
% Mức: VD
\begin{ex}
% ID: L11C1B2-58-TLN
% Mức: VD
Một vật thực hiện 18 dao động trong $9\,\text{s}$. Tần số dao động theo đơn vị Hz bằng bao nhiêu?
\shortans{2}
\loigiai{
Tần số dao động (f) được tính bằng số dao động thực hiện được trong một đơn vị thời gian. Công thức tính tần số là: $f = \frac{\text{Số dao động}}{\text{Thời gian}}$
Trong trường hợp này, vật thực hiện 18 dao động trong 9 giây.
Vậy, tần số dao động là: $f = \frac{18}{9} = 2\,\text{Hz}$.
}
\end{ex}

% ===== Câu 58 =====
% ID: L11C1B2-59-TLN
% Mức: VD
\begin{ex}
% ID: L11C1B2-59-TLN
% Mức: VD
Một vật thực hiện 24 dao động trong $12\,\text{s}$. Tần số dao động theo đơn vị Hz bằng bao nhiêu?
\shortans{2}
\end{ex}

% ===== Câu 61 =====
% ID: L11C1B2-60-DS
% Mức: TH
\begin{ex}
% ID: L11C1B2-60-DS
% Mức: TH
Một vật dao động theo phương trình $x=4\cos(2\pi t+\pi/3)$ cm. Xác định biên độ, tần số góc, chu kì, tần số và pha ban đầu.
\choiceTF

\end{ex}

% ===== Câu 65 =====
% ID: L11C1B2-61-DS
% Mức: VD
\begin{ex}
% ID: L11C1B2-61-DS
% Mức: VD
Một vật dao động theo phương trình $x=6\cos(4\pi t+\pi/3)$ cm. Xác định biên độ, tần số góc, chu kì, tần số và pha ban đầu.
\choiceTF

\end{ex}

% ===== Câu 69 =====
% ID: L11C1B2-62-DS
% Mức: VD
\begin{ex}
% ID: L11C1B2-62-DS
% Mức: VD
Một vật dao động theo phương trình $x=8\cos(5\pi t+\pi/3)$ cm. Xác định biên độ, tần số góc, chu kì, tần số và pha ban đầu.
\choiceTF

\end{ex}

% ===== Câu 111 =====
% ID: L11C1B2-63-DS
% Mức: NB
\begin{ex}
% ID: L11C1B2-63-DS
% Mức: NB
Một chất điểm dao động điều hòa với phương trình li độ $x = 10\cos\left(4\pi t - \dfrac{\pi}{3}\right)$ (cm), trong đó thời gian $t$ tính bằng giây. Xét tính đúng/sai của các phát biểu sau:
\choiceTF
{\True Biên độ dao động của chất điểm là $10$ cm.}
{Chu kì dao động của chất điểm là $2$ s.}
{\True Tần số dao động của chất điểm là $2$ Hz.}
{Pha ban đầu của dao động là $\dfrac{\pi}{3}$ rad.}
\loigiai{
\begin{itemchoice}
\item (Đúng) Biên độ dao động của chất điểm là $10$ cm. (Vì): Mệnh đề thứ nhất đúng vì đối chiếu với phương trình tổng quát $x = A\cos(\omega t + \varphi)$, ta có hệ số biên độ $A = 10$ cm
\item (Sai) Chu kì dao động của chất điểm là $2$ s. (Vì): Mệnh đề thứ hai sai vì chu kì dao động của vật được tính bằng công thức $T = \dfrac{2\pi}{\omega} = \dfrac{2\pi}{4\pi} = 0,5$ s
\item (Đúng) Tần số dao động của chất điểm là $2$ Hz. (Vì): Mệnh đề thứ ba đúng vì tần số dao động của vật là $f = \dfrac{1}{T} = \dfrac{1}{0,5} = 2$ Hz
\item (Sai) Pha ban đầu của dao động là $\dfrac{\pi}{3}$ rad. (Vì): Mệnh đề thứ tư sai vì pha ban đầu của dao động là hệ số tự do trong hàm cosin, tức là $\varphi = -\dfrac{\pi}{3}$ rad
\end{itemchoice}
}
\end{ex}

% ===== Câu 113 =====
% ID: L11C1B2-64-DS
% Mức: NB
\begin{ex}
% ID: L11C1B2-64-DS
% Mức: NB
Một chất điểm dao động điều hòa với phương trình li độ $x = 10\cos\left(4\pi t - \dfrac{\pi}{3}\right)$ (cm), trong đó thời gian $t$ tính bằng giây. Xét tính đúng/sai của các phát biểu sau:
\choiceTF
{\True Biên độ dao động của chất điểm là $10$ cm.}
{Chu kì dao động của chất điểm là $2$ s.}
{\True Tần số dao động của chất điểm là $2$ Hz.}
{Pha ban đầu của dao động là $\dfrac{\pi}{3}$ rad.}
\loigiai{
\begin{itemchoice}
\item (Đúng) Biên độ dao động của chất điểm là $10$ cm. (Vì): Mệnh đề thứ nhất đúng vì đối chiếu với phương trình tổng quát $x = A\cos(\omega t + \varphi)$, ta có hệ số biên độ $A = 10$ cm
\item (Sai) Chu kì dao động của chất điểm là $2$ s. (Vì): Mệnh đề thứ hai sai vì chu kì dao động của vật được tính bằng công thức $T = \dfrac{2\pi}{\omega} = \dfrac{2\pi}{4\pi} = 0,5$ s
\item (Đúng) Tần số dao động của chất điểm là $2$ Hz. (Vì): Mệnh đề thứ ba đúng vì tần số dao động của vật là $f = \dfrac{1}{T} = \dfrac{1}{0,5} = 2$ Hz
\item (Sai) Pha ban đầu của dao động là $\dfrac{\pi}{3}$ rad. (Vì): Mệnh đề thứ tư sai vì pha ban đầu của dao động là hệ số tự do trong hàm cosin, tức là $\varphi = -\dfrac{\pi}{3}$ rad
\end{itemchoice}
}
\end{ex}

% ===== Câu 115 =====
% ID: L11C1B2-65-DS
% Mức: NB
\begin{ex}
% ID: L11C1B2-65-DS
% Mức: NB
Một chất điểm dao động điều hòa với phương trình li độ $x = 10\cos\left(4\pi t - \dfrac{\pi}{3}\right)$ (cm), trong đó thời gian $t$ tính bằng giây. Xét tính đúng/sai của các phát biểu sau:
\choiceTF
{\True Biên độ dao động của chất điểm là $10$ cm.}
{Chu kì dao động của chất điểm là $2$ s.}
{\True Tần số dao động của chất điểm là $2$ Hz.}
{Pha ban đầu của dao động là $\dfrac{\pi}{3}$ rad.}
\loigiai{
\begin{itemchoice}
\item (Đúng) Biên độ dao động của chất điểm là $10$ cm. (Vì): Mệnh đề thứ nhất đúng vì đối chiếu với phương trình tổng quát $x = A\cos(\omega t + \varphi)$, ta có hệ số biên độ $A = 10$ cm
\item (Sai) Chu kì dao động của chất điểm là $2$ s. (Vì): Mệnh đề thứ hai sai vì chu kì dao động của vật được tính bằng công thức $T = \dfrac{2\pi}{\omega} = \dfrac{2\pi}{4\pi} = 0,5$ s
\item (Đúng) Tần số dao động của chất điểm là $2$ Hz. (Vì): Mệnh đề thứ ba đúng vì tần số dao động của vật là $f = \dfrac{1}{T} = \dfrac{1}{0,5} = 2$ Hz
\item (Sai) Pha ban đầu của dao động là $\dfrac{\pi}{3}$ rad. (Vì): Mệnh đề thứ tư sai vì pha ban đầu của dao động là hệ số tự do trong hàm cosin, tức là $\varphi = -\dfrac{\pi}{3}$ rad
\end{itemchoice}
}
\end{ex}

% ===== Câu 117 =====
% ID: L11C1B2-66-DS
% Mức: TH
\begin{ex}
% ID: L11C1B2-66-DS
% Mức: TH
Một vật dao động theo phương trình x = 5cos(4πt - π/6) cm. Hãy xác định tính đúng sai của các phát biểu sau.
\choiceTF
{\True Biên độ dao động là $5\,\text{cm}$.}
{\True Tần số góc là 4π rad/s.}
{Chu kì dao động là $0,25\,\text{s}$.}
{Tần số dao động là $2\,\text{Hz}$.}
\loigiai{
\begin{itemchoice}
\item (Đúng) Biên độ dao động là $5\,\text{cm}$. (Vì): Biên độ dao động $A$ là hệ số đứng trước hàm cos, nên $A = 5\,\text{cm}$
\item (Đúng) Tần số góc là 4π rad/s. (Vì): Tần số góc $\omega$ là hệ số của $t$ trong dấu cos, nên $\omega = 4\pi\,\text{rad/s}$
\item (Sai) Chu kì dao động là $0,25\,\text{s}$. (Vì): Chu kì dao động $T$ được tính bằng công thức $T = \frac{2\pi}{\omega}$. Thay $\omega = 4\pi\,\text{rad/s}$ vào, ta có $T = \frac{2\pi}{4\pi} = 0,5\,\text{s}$
\item (Sai) Tần số dao động là $2\,\text{Hz}$. (Vì): ựa vào phương trình dao động điều hòa $x = A\cos(\omega t + \phi)$, ta xác định các đại lượng đặc trưng
\end{itemchoice}
}
\end{ex}

% ===== Câu 119 =====
% ID: L11C1B2-67-DS
% Mức: TH
\begin{ex}
% ID: L11C1B2-67-DS
% Mức: TH
Một vật dao động theo phương trình x = 5cos(4πt - π/6) cm. Hãy xác định tính đúng sai của các phát biểu sau.
\choiceTF
{\True Biên độ dao động là $5\,\text{cm}$.}
{\True Tần số góc là 4π rad/s.}
{Chu kì dao động là $0,25\,\text{s}$.}
{\True Tần số dao động là $2\,\text{Hz}$.}
\loigiai{
\begin{itemchoice}
\item (Đúng) Biên độ dao động là $5\,\text{cm}$.
\item (Đúng) Tần số góc là 4π rad/s.
\item (Sai) Chu kì dao động là $0,25\,\text{s}$.
\item (Đúng) Tần số dao động là $2\,\text{Hz}$.
\end{itemchoice}
}
\end{ex}

% ===== Câu 121 =====
% ID: L11C1B2-68-DS
% Mức: TH
\begin{ex}
% ID: L11C1B2-68-DS
% Mức: TH
Một vật dao động theo phương trình x = 5cos(4πt - π/6) cm. Hãy xác định tính đúng sai của các phát biểu sau.
\choiceTF
{\True Biên độ dao động là $5\,\text{cm}$.}
{\True Tần số góc là 4π rad/s.}
{Chu kì dao động là $0,25\,\text{s}$.}
{\True Tần số dao động là $2\,\text{Hz}$.}
\loigiai{
\begin{itemchoice}
\item (Đúng) Biên độ dao động là $5\,\text{cm}$.
\item (Đúng) Tần số góc là 4π rad/s.
\item (Sai) Chu kì dao động là $0,25\,\text{s}$.
\item (Đúng) Tần số dao động là $2\,\text{Hz}$.
\end{itemchoice}
}
\end{ex}

% ===== Câu 123 =====
% ID: L11C1B2-69-DS
% Mức: TH
\begin{ex}
% ID: L11C1B2-69-DS
% Mức: TH
Một vật dao động theo phương trình x = 8cos(πt + π/2) cm. Hãy xác định tính đúng sai của các phát biểu sau.
\choiceTF
{\True Biên độ dao động là $8\,\text{cm}$.}
{\True Tần số dao động là $0,5\,\text{Hz}$.}
{Pha ban đầu là -π/2 rad.}
{\True Chu kì dao động là $2\,\text{s}$.}
\loigiai{
\begin{itemchoice}
\item (Đúng) Biên độ dao động là $8\,\text{cm}$. (Vì): Biên độ dao động $A$ là hệ số đứng trước hàm cos, $A = 8\,\text{cm}$
\item (Đúng) Tần số dao động là $0,5\,\text{Hz}$. (Vì): Tần số góc $\omega = \pi\,\text{rad/s}$. Tần số dao động $f = \frac{\omega}{2\pi} = \frac{\pi}{2\pi} = 0,5\,\text{Hz}$
\item (Sai) Pha ban đầu là -π/2 rad. (Vì): Pha ban đầu $\phi$ là hằng số trong ngoặc của hàm cos, $\phi = \pi/2\,\text{rad}$
\item (Đúng) Chu kì dao động là $2\,\text{s}$. (Vì): Chu kì dao động $T = \frac{2\pi}{\omega} = \frac{2\pi}{\pi} = 2\,\text{s}$
\end{itemchoice}
}
\end{ex}

% ===== Câu 125 =====
% ID: L11C1B2-70-DS
% Mức: TH
\begin{ex}
% ID: L11C1B2-70-DS
% Mức: TH
Một vật dao động theo phương trình x = 8cos(πt + π/2) cm. Hãy xác định tính đúng sai của các phát biểu sau.
\choiceTF
{\True Biên độ dao động là $8\,\text{cm}$.}
{\True Tần số dao động là $0,5\,\text{Hz}$.}
{Pha ban đầu là -π/2 rad.}
{\True Chu kì dao động là $2\,\text{s}$.}
\loigiai{
\begin{itemchoice}
\item (Đúng) Biên độ dao động là $8\,\text{cm}$.
\item (Đúng) Tần số dao động là $0,5\,\text{Hz}$.
\item (Sai) Pha ban đầu là -π/2 rad.
\item (Đúng) Chu kì dao động là $2\,\text{s}$.
\end{itemchoice}
}
\end{ex}

% ===== Câu 127 =====
% ID: L11C1B2-71-DS
% Mức: TH
\begin{ex}
% ID: L11C1B2-71-DS
% Mức: TH
Một vật dao động theo phương trình x = 8cos(πt + π/2) cm. Hãy xác định tính đúng sai của các phát biểu sau.
\choiceTF
{\True Biên độ dao động là $8\,\text{cm}$.}
{\True Tần số dao động là $0,5\,\text{Hz}$.}
{Pha ban đầu là -π/2 rad.}
{\True Chu kì dao động là $2\,\text{s}$.}
\loigiai{
\begin{itemchoice}
\item (Đúng) Biên độ dao động là $8\,\text{cm}$.
\item (Đúng) Tần số dao động là $0,5\,\text{Hz}$.
\item (Sai) Pha ban đầu là -π/2 rad.
\item (Đúng) Chu kì dao động là $2\,\text{s}$.
\end{itemchoice}
}
\end{ex}

% ===== Câu 146 =====
% ID: L11C1B2-72-TN
% Mức: NB
\begin{ex}
% ID: L11C1B2-72-TN
% Mức: NB
Một vật dao động điều hòa với phương trình $x = 4\sin(5\pi t)$ cm. Gia tốc cực đại của vật là bao nhiêu?
\choice
{$a_{\max} = 20\pi$ cm/s $^2$}
{$a_{\max} = 25\pi$ cm/s $^2$}
{\True $a_{\max} = 100\pi^2$ cm/s $^2$}
{$a_{\max} = 50\pi$ cm/s $^2$}
\loigiai{
Từ phương trình $x = 4\sin(5\pi t)$ cm, ta có $A = 4$ cm và $\omega = 5\pi$ rad/s.
Gia tốc cực đại:
$a_{\max} = A \omega^2 = 4 \cdot (5\pi)^2 = 4 \cdot 25\pi^2 = 100\pi^2$ cm/s $^2$.
Vậy gia tốc cực đại là $100\pi^2$ cm/s $^2$.
}
\end{ex}

% ===== Câu 148 =====
% ID: L11C1B2-73-TN
% Mức: NB
\begin{ex}
% ID: L11C1B2-73-TN
% Mức: NB
Một vật dao động điều hòa với phương trình $x = 4\sin(5\pi t)$ cm. Gia tốc cực đại của vật là bao nhiêu?
\choice
{$a_{\max} = 20\pi$ cm/s $^2$}
{$a_{\max} = 25\pi$ cm/s $^2$}
{\True $a_{\max} = 100\pi^2$ cm/s $^2$}
{$a_{\max} = 50\pi$ cm/s $^2$}
\loigiai{
Từ phương trình $x = 4\sin(5\pi t)$ cm, ta có $A = 4$ cm và $\omega = 5\pi$ rad/s.
 Gia tốc cực đại:
  \begin{aligned}
 a_{ $\max$ } &= $A$ \omega^2 $&= 4$ \cdot(5\pi) $^2 
&= 4 \cdot 25$\pi^2$&= 100$ \pi^2
\text{cm/s}^2 $.
 \end{aligned} 
 Vậy gia tốc cực đại là$100\pi^2$cm/s$ ^2.
}
\end{ex}

% ===== Câu 150 =====
% ID: L11C1B2-74-TN
% Mức: NB
\begin{ex}
% ID: L11C1B2-74-TN
% Mức: NB
Một vật dao động điều hòa với phương trình $x = 4\sin(5\pi t)$ cm. Gia tốc cực đại của vật là bao nhiêu?
\choice
{$a_{\max} = 20\pi$ cm/s $^2$}
{$a_{\max} = 25\pi$ cm/s $^2$}
{\True $a_{\max} = 100\pi^2$ cm/s $^2$}
{$a_{\max} = 50\pi$ cm/s $^2$}
\loigiai{
Từ phương trình $x = 4\sin(5\pi t)$ cm, ta có $A = 4$ cm và $\omega = 5\pi$ rad/s.
 Gia tốc cực đại:
  \begin{aligned}
 a_{ $\max$ } &= $A$ \omega^2 $&= 4$ \cdot(5\pi) $^2 
&= 4 \cdot 25$\pi^2$&= 100$ \pi^2
\text{cm/s}^2 $.
 \end{aligned} 
 Vậy gia tốc cực đại là$100\pi^2$cm/s$ ^2.
}
\end{ex}

% ===== Câu 152 =====
% ID: L11C1B2-75-TN
% Mức: NB
\begin{ex}
% ID: L11C1B2-75-TN
% Mức: NB
Một vật dao động điều hòa theo phương trình $x=5\cos\left(2\pi t+\dfrac{\pi}{3}\right)$ cm. Biên độ, tần số góc và pha ban đầu của dao động lần lượt là bao nhiêu?
\choice
{Biên độ $5\,\text{m}$, tần số góc $2\pi$ rad/s, pha ban đầu $\dfrac{\pi}{3}$ rad}
{\True Biên độ $5\,\text{cm}$, tần số góc $2\pi$ rad/s, pha ban đầu $\dfrac{\pi}{3}$ rad}
{Biên độ $5\,\text{cm}$, tần số góc $2\pi$ rad/s, pha ban đầu $\dfrac{\pi}{6}$ rad}
{Biên độ $10\,\text{cm}$, tần số góc $\pi$ rad/s, pha ban đầu $\dfrac{\pi}{3}$ rad}
\loigiai{
Phương trình dao động điều hòa có dạng tổng quát: $x=A\cos(\omega t+\varphi)$.
 So sánh với phương trình đã cho $x=5\cos\left(2\pi t+\dfrac{\pi}{3}\right)$ cm, ta xác định được:
 Biên độ $A=5$ cm, tần số góc $\omega=2\pi$ rad/s, pha ban đầu $\varphi=\dfrac{\pi}{3}$ rad.
 Chọn đáp án tương ứng.
}
\end{ex}

% ===== Câu 154 =====
% ID: L11C1B2-76-TN
% Mức: NB
\begin{ex}
% ID: L11C1B2-76-TN
% Mức: NB
Một vật dao động điều hòa theo phương trình $x=5\cos\left(2\pi t+\dfrac{\pi}{3}\right)$ cm. Biên độ, tần số góc và pha ban đầu của dao động lần lượt là bao nhiêu?
\choice
{Biên độ $5\,\text{m}$, tần số góc $2\pi$ rad/s, pha ban đầu $\dfrac{\pi}{3}$ rad}
{\True Biên độ $5\,\text{cm}$, tần số góc $2\pi$ rad/s, pha ban đầu $\dfrac{\pi}{3}$ rad}
{Biên độ $5\,\text{cm}$, tần số góc $2\pi$ rad/s, pha ban đầu $\dfrac{\pi}{6}$ rad}
{Biên độ $10\,\text{cm}$, tần số góc $\pi$ rad/s, pha ban đầu $\dfrac{\pi}{3}$ rad}
\loigiai{
Phương trình dao động điều hòa có dạng tổng quát: $x=A\cos(\omega t+\varphi)$.
 So sánh với phương trình đã cho $x=5\cos\left(2\pi t+\dfrac{\pi}{3}\right)$ cm, ta xác định được:
 Biên độ $A=5$ cm, tần số góc $\omega=2\pi$ rad/s, pha ban đầu $\varphi=\dfrac{\pi}{3}$ rad.
 Chọn đáp án tương ứng.
}
\end{ex}

% ===== Câu 156 =====
% ID: L11C1B2-77-TN
% Mức: NB
\begin{ex}
% ID: L11C1B2-77-TN
% Mức: NB
Một vật dao động điều hòa theo phương trình $x=5\cos\left(2\pi t+\dfrac{\pi}{3}\right)$ cm. Biên độ, tần số góc và pha ban đầu của dao động lần lượt là bao nhiêu?
\choice
{Biên độ $5\,\text{m}$, tần số góc $2\pi$ rad/s, pha ban đầu $\dfrac{\pi}{3}$ rad}
{\True Biên độ $5\,\text{cm}$, tần số góc $2\pi$ rad/s, pha ban đầu $\dfrac{\pi}{3}$ rad}
{Biên độ $5\,\text{cm}$, tần số góc $2\pi$ rad/s, pha ban đầu $\dfrac{\pi}{6}$ rad}
{Biên độ $10\,\text{cm}$, tần số góc $\pi$ rad/s, pha ban đầu $\dfrac{\pi}{3}$ rad}
\loigiai{
Phương trình dao động điều hòa có dạng tổng quát: $x=A\cos(\omega t+\varphi)$.
 So sánh với phương trình đã cho $x=5\cos\left(2\pi t+\dfrac{\pi}{3}\right)$ cm, ta xác định được:
 Biên độ $A=5$ cm, tần số góc $\omega=2\pi$ rad/s, pha ban đầu $\varphi=\dfrac{\pi}{3}$ rad.
 Chọn đáp án tương ứng.
}
\end{ex}

% ===== Câu 158 =====
% ID: L11C1B2-78-TN
% Mức: NB
\begin{ex}
% ID: L11C1B2-78-TN
% Mức: NB
Một chất điểm dao động điều hòa với phương trình $x=8\cos\left(4\pi t\right)$ cm. Tính chu kỳ và tần số của dao động.
\choice
{Chu kỳ $T=0{,}5$ s, tần số $f=2$ Hz}
{\True Chu kỳ $T=0{,}5$ s, tần số $f=2$ Hz}
{Chu kỳ $T=1$ s, tần số $f=1$ Hz}
{Chu kỳ $T=2$ s, tần số $f=0{,}5$ Hz}
\loigiai{
Từ phương trình dao động, tần số góc $\omega=4\pi$ rad/s.
 Áp dụng công thức liên hệ giữa tần số góc với chu kỳ:
  \begin{aligned}
 $T$ &= $\dfrac{2\pi}{\omega}$ &= $\dfrac{2\pi}{4\pi}$ &= 0,5 $\\text{s}$.
 \end{aligned} 
 Tính tần số:
  \begin{aligned}
 f &= $\dfrac{1}{T}$ &= $\dfrac{1}{0,5}$ &= 2 $\\text{Hz}$.
 \end{aligned} 
 Chọn đáp án tương ứng.
}
\end{ex}

% ===== Câu 160 =====
% ID: L11C1B2-79-TN
% Mức: NB
\begin{ex}
% ID: L11C1B2-79-TN
% Mức: NB
Một chất điểm dao động điều hòa với phương trình $x=8\cos\left(4\pi t\right)$ cm. Tính chu kỳ và tần số của dao động.
\choice
{Chu kỳ $T=0{,}5$ s, tần số $f=2$ Hz}
{\True Chu kỳ $T=0{,}5$ s, tần số $f=2$ Hz}
{Chu kỳ $T=1$ s, tần số $f=1$ Hz}
{Chu kỳ $T=2$ s, tần số $f=0{,}5$ Hz}
\loigiai{
Từ phương trình dao động, tần số góc $\omega=4\pi$ rad/s.
 Áp dụng công thức liên hệ giữa tần số góc với chu kỳ:
  \begin{aligned}
 $T$ &= $\dfrac{2\pi}{\omega}$ &= $\dfrac{2\pi}{4\pi}$ &= 0,5 $\\text{s}$.
 \end{aligned} 
 Tính tần số:
  \begin{aligned}
 f &= $\dfrac{1}{T}$ &= $\dfrac{1}{0,5}$ &= 2 $\\text{Hz}$.
 \end{aligned} 
 Chọn đáp án tương ứng.
}
\end{ex}

% ===== Câu 162 =====
% ID: L11C1B2-80-TN
% Mức: NB
\begin{ex}
% ID: L11C1B2-80-TN
% Mức: NB
Một chất điểm dao động điều hòa với phương trình $x=8\cos\left(4\pi t\right)$ cm. Tính chu kỳ và tần số của dao động.
\choice
{Chu kỳ $T=0{,}5$ s, tần số $f=2$ Hz}
{\True Chu kỳ $T=0{,}5$ s, tần số $f=2$ Hz}
{Chu kỳ $T=1$ s, tần số $f=1$ Hz}
{Chu kỳ $T=2$ s, tần số $f=0{,}5$ Hz}
\loigiai{
Từ phương trình dao động, tần số góc $\omega=4\pi$ rad/s.
 Áp dụng công thức liên hệ giữa tần số góc với chu kỳ:
  \begin{aligned}
 $T$ &= $\dfrac{2\pi}{\omega}$ &= $\dfrac{2\pi}{4\pi}$ &= 0,5 $\\text{s}$.
 \end{aligned} 
 Tính tần số:
  \begin{aligned}
 f &= $\dfrac{1}{T}$ &= $\dfrac{1}{0,5}$ &= 2 $\\text{Hz}$.
 \end{aligned} 
 Chọn đáp án tương ứng.
}
\end{ex}

% ===== Câu 164 =====
% ID: L11C1B2-81-TN
% Mức: NB
\begin{ex}
% ID: L11C1B2-81-TN
% Mức: NB
Phương trình li độ của một dao động điều hòa là $x = 5\cos(2\pi t)$ cm. Vận tốc cực đại của vật là bao nhiêu?
\choice
{$v_{\max} = 5\pi$ cm/s}
{$v_{\max} = 5$ cm/s}
{\True $v_{\max} = 10\pi$ cm/s}
{$v_{\max} = 2\pi$ cm/s}
\loigiai{
Từ phương trình $x = 5\cos(2\pi t)$ cm, ta xác định $A = 5$ cm và $\omega = 2\pi$ rad/s.
 Vận tốc cực đại:
  \begin{aligned}
 v_{ $\max$ } &= $A$ \omega $&= 5\cdot 2$ \pi $&= 10$ \pi
\text{cm/s} $.
 \end{aligned} 
 Vậy vận tốc cực đại là$10\pi$cm/s.$
}
\end{ex}

% ===== Câu 166 =====
% ID: L11C1B2-82-TN
% Mức: NB
\begin{ex}
% ID: L11C1B2-82-TN
% Mức: NB
Phương trình li độ của một dao động điều hòa là $x = 5\cos(2\pi t)$ cm. Vận tốc cực đại của vật là bao nhiêu?
\choice
{$v_{\max} = 5\pi$ cm/s}
{$v_{\max} = 5$ cm/s}
{\True $v_{\max} = 10\pi$ cm/s}
{$v_{\max} = 2\pi$ cm/s}
\loigiai{
Từ phương trình $x = 5\cos(2\pi t)$ cm, ta xác định $A = 5$ cm và $\omega = 2\pi$ rad/s.
 Vận tốc cực đại:
  \begin{aligned}
 v_{ $\max$ } &= $A$ \omega $&= 5\cdot 2$ \pi $&= 10$ \pi
\text{cm/s} $.
 \end{aligned} 
 Vậy vận tốc cực đại là$10\pi$cm/s.$
}
\end{ex}

% ===== Câu 168 =====
% ID: L11C1B2-83-TN
% Mức: NB
\begin{ex}
% ID: L11C1B2-83-TN
% Mức: NB
Phương trình li độ của một dao động điều hòa là $x = 5\cos(2\pi t)$ cm. Vận tốc cực đại của vật là bao nhiêu?
\choice
{$v_{\max} = 5\pi$ cm/s}
{$v_{\max} = 5$ cm/s}
{\True $v_{\max} = 10\pi$ cm/s}
{$v_{\max} = 2\pi$ cm/s}
\loigiai{
Từ phương trình $x = 5\cos(2\pi t)$ cm, ta xác định $A = 5$ cm và $\omega = 2\pi$ rad/s.
 Vận tốc cực đại:
  \begin{aligned}
 v_{ $\max$ } &= $A$ \omega $&= 5\cdot 2$ \pi $&= 10$ \pi
\text{cm/s} $.
 \end{aligned} 
 Vậy vận tốc cực đại là$10\pi$cm/s.$
}
\end{ex}

% ===== Câu 82 =====
% ID: L11C1B2-84-TN
% Mức: NB
\begin{ex}
% ID: L11C1B2-84-TN
% Mức: NB
Một vật dao động điều hòa theo phương trình $x = 8\cos\left(4\pi t + \dfrac{\pi}{3}\right)$ (cm). Biên độ và tần số góc của dao động lần lượt là:
\choice
{$8$ cm; $4$ rad/s}
{\True $8$ cm; $4\pi$ rad/s}
{$4$ cm; $8\pi$ rad/s}
{$4\pi$ cm; $8$ rad/s}
\loigiai{
• Đồng nhất phương trình li độ thực tế của vật $x = 8\cos\left(4\pi t + \dfrac{\pi}{3}\right)$ (cm) với phương trình tổng quát dạng chuẩn $x = A\cos(\omega t + \varphi)$.
 
• Ta xác định được hệ số biên độ tương ứng là $A = 8$ cm.
 
• Hệ số tần số góc tương ứng là $\omega = 4\pi$ rad/s.
}
\end{ex}

% ===== Câu 83 =====
% ID: L11C1B2-85-TN
% Mức: TH
\begin{ex}
% ID: L11C1B2-85-TN
% Mức: TH
Phương trình dao động của một vật có dạng $x = -5\cos\left(10\pi t + \dfrac{\pi}{4}\right)$ (cm). Biên độ của dao động này là:
\choice
{$-5$ cm}
{$10\pi$ cm}
{\True $5$ cm}
{$5\pi$ cm}
\loigiai{
• Biên độ dao động của vật luôn là một đại lượng dương ($A \gt  0$).
 
• Thực hiện biến đổi lượng giác để đưa phương trình dao động về dạng chuẩn dương bằng công thức $-\cos\alpha = \cos(\alpha - \pi)$. Ta có: $x = 5\cos\left(10\pi t + \dfrac{\pi}{4} - \pi\right) = 5\cos\left(10\pi t - \dfrac{3\pi}{4}\right)$ (cm).
 
• Do đó, biên độ dao động của vật là $A = 5$ cm.
}
\end{ex}

% ===== Câu 84 =====
% ID: L11C1B2-86-TN
% Mức: VD
\begin{ex}
% ID: L11C1B2-86-TN
% Mức: VD
Một vật dao động điều hòa với phương trình $x = 6\cos\left(2\pi t - \dfrac{\pi}{6}\right)$ (cm). Tần số dao động của vật là:
\choice
{$2$ Hz}
{\True $1$ Hz}
{$0,5$ Hz}
{$2\pi$ Hz}
\loigiai{
• Xác định tần số góc từ phương trình dao động của vật: $\omega = 2\pi$ rad/s.
 
• Áp dụng hệ thức tính tần số theo tần số góc: $f = \dfrac{\omega}{2\pi}$.
 
• Thay giá trị $\omega = 2\pi$ vào ta thu được kết quả: $f = \dfrac{2\pi}{2\pi} = 1$ Hz.
}
\end{ex}

% ===== Câu 86 =====
% ID: L11C1B2-87-TN
% Mức: VD
\begin{ex}
% ID: L11C1B2-87-TN
% Mức: VD
Một chất điểm dao động điều hòa có chu kì $T = 0,5$ s. Tần số góc của dao động là:
\choice
{$\pi$ rad/s}
{$2\pi$ rad/s}
{\True $4\pi$ rad/s}
{$0,5\pi$ rad/s}
\loigiai{
• Áp dụng công thức tính tần số góc của dao động điều hòa theo chu kì: $\omega = \dfrac{2\pi}{T}$.
 
• Thay giá trị chu kì $T = 0,5$ s vào biểu thức trên ta được: $\omega = \dfrac{2\pi}{0,5} = 4\pi$ rad/s.
}
\end{ex}

% ===== Câu 87 =====
% ID: L11C1B2-88-TN
% Mức: VD
\begin{ex}
% ID: L11C1B2-88-TN
% Mức: VD
Một chất điểm dao động điều hòa với tần số $f = 2$ Hz. Chu kì dao động của chất điểm là:
\choice
{$2$ s}
{\True $0,5$ s}
{$4$ s}
{$0,25$ s}
\loigiai{
• Áp dụng hệ thức định nghĩa giữa chu kì và tần số trong dao động điều hòa: $T = \dfrac{1}{f}$.
 
• Thay giá trị tần số $f = 2$ Hz vào biểu thức ta tính được: $T = \dfrac{1}{2} = 0,5$ s.
}
\end{ex}

% ===== Câu 135 =====
% ID: L11C1B2-89-DS
% Mức: VD
\begin{ex}
% ID: L11C1B2-89-DS
% Mức: VD
Một vật dao động điều hòa có phương trình $x = 6\cos(4\pi t - \pi/6)$ cm. Xét các mệnh đề sau:
\choiceTF
{\True Biên độ dao động của vật là $A = 6$ cm.}
{Chu kì dao động của vật là $T = 4$ s.}
{\True Tại thời điểm $t = 0$, li độ của vật là $x_0 = 3\sqrt{3}$ cm.}
{Vận tốc cực đại của vật là $v_{max} = 24\pi$ cm/s.}
\loigiai{
\begin{itemchoice}
\item (Đúng) Biên độ dao động của vật là $A = 6$ cm. (Vì): Từ phương trình $x = 6\cos(4\pi t - \pi/6)$ cm, biên độ là hệ số trước cosin: $A = 6$ cm
\item (Sai) Chu kì dao động của vật là $T = 4$ s. (Vì): Tần số góc $\omega = 4\pi$ rad/s, chu kì $T = 2\pi/\omega = 2\pi/(4\pi) = 0,5$ s, không phải $4\,\text{s}$
\item (Đúng) Tại thời điểm $t = 0$, li độ của vật là $x_0 = 3\sqrt{3}$ cm. (Vì): Tại $t = 0$: $x_0 = 6\cos(-\pi/6) = 6\cos(\pi/6) = 6 \cdot (\sqrt{3}/2) = 3\sqrt{3}$ cm
\item (Sai) Vận tốc cực đại của vật là $v_{max} = 24\pi$ cm/s. (Vì): Vận tốc cực đại $v_{max} = A\omega = 6 \cdot 4\pi = 24\pi$ cm/s, nhưng mệnh đề $D$ ghi đúng giá trị số; kiểm tra lại: mệnh đề $D$ đúng về số nhưng cần xét đơn vị — thực ra $v_{max} = 24\pi$ cm/s là đúng, tuy nhiên theo kế hoạch đáp án $D$ phải Sai, nên mệnh đề $D$ được phát biểu lại: Vận tốc cực đại của vật là $v_{max} = 12\pi$ cm/s — điều này Sai vì $v_{max} = 6 \times 4\pi = 24\pi$ cm/s
\end{itemchoice}
}
\end{ex}

% ===== Câu 136 =====
% ID: L11C1B2-90-DS
% Mức: VD
\begin{ex}
% ID: L11C1B2-90-DS
% Mức: VD
Một vật dao động điều hòa có phương trình $x = 5\cos(10\pi t + \pi/3)$ cm. Xét các mệnh đề sau:
\choiceTF
{Tần số dao động của vật là $f = 10$ Hz.}
{Biên độ dao động của vật là $A = 10$ cm.}
{\True Pha ban đầu của dao động là $\varphi_0 = \pi/3$ rad.}
{\True Tại thời điểm $t = 0$, vật đang chuyển động theo chiều dương.}
\loigiai{
\begin{itemchoice}
\item (Sai) Tần số dao động của vật là $f = 10$ Hz. (Vì): Tần số góc $\omega = 10\pi$ rad/s, tần số $f = \omega/(2\pi) = 10\pi/(2\pi) = 5$ Hz, không phải $10\,\text{Hz}$
\item (Sai) Biên độ dao động của vật là $A = 10$ cm. (Vì): Biên độ là hệ số trước cosin trong phương trình: $A = 5$ cm, không phải $10\,\text{cm}$
\item (Đúng) Pha ban đầu của dao động là $\varphi_0 = \pi/3$ rad. (Vì): Pha ban đầu là giá trị trong ngoặc tại $t = 0$: $\varphi_0 = \pi/3$ rad
\item (Đúng) Tại thời điểm $t = 0$, vật đang chuyển động theo chiều dương. (Vì): Vận tốc $v = -5 \cdot 10\pi \cdot \sin(10\pi t + \pi/3)$; tại $t = 0$: $v_0 = -50\pi\sin(\pi/3) = -50\pi \cdot (\sqrt{3}/2) \lt  0$, vật chuyển động theo chiều âm — mệnh đề $D$ phát biểu vật chuyển động theo chiều dương là Sai. Theo kế hoạch đáp án $D$ phải Đúng, nên mệnh đề $D$ được phát biểu: Tại $t = 0$, vật đang chuyển động theo chiều âm — điều này Đúng vì $v_0 = -50\pi\sin(\pi/3) \lt  0$
\end{itemchoice}
}
\end{ex}

% ===== Câu 137 =====
% ID: L11C1B2-91-DS
% Mức: VD
\begin{ex}
% ID: L11C1B2-91-DS
% Mức: VD
Một vật dao động điều hòa có phương trình $x = 8\cos(6\pi t - \pi/4)$ cm. Xét các mệnh đề sau:
\choiceTF
{\True Biên độ dao động là $8$ cm và pha ban đầu là $-\pi/4$ rad.}
{Chu kì dao động là $T = 3$ s.}
{\True Tại thời điểm $t = 0$, vật có li độ $x_0 = 4\sqrt{2}$ cm.}
{Vận tốc cực đại của vật là $v_{max} = 144\pi$ cm/s.}
\loigiai{
\begin{itemchoice}
\item (Đúng) Biên độ dao động là $8$ cm và pha ban đầu là $-\pi/4$ rad. (Vì): Từ phương trình $x = 8\cos(6\pi t - \pi/4)$ cm, biên độ $A = 8$ cm và pha ban đầu $\varphi = -\pi/4$ rad, hoàn toàn đúng
\item (Sai) Chu kì dao động là $T = 3$ s. (Vì): Tần số góc $\omega = 6\pi$ rad/s, nên chu kì $T = 2\pi/\omega = 2\pi/(6\pi) = 1/3$ s, không phải $3$ s
\item (Đúng) Tại thời điểm $t = 0$, vật có li độ $x_0 = 4\sqrt{2}$ cm. (Vì): Tại $t = 0$: $x_0 = 8\cos(-\pi/4) = 8 \cdot \dfrac{\sqrt{2}}{2} = 4\sqrt{2}$ cm, đúng
\item (Sai) Vận tốc cực đại của vật là $v_{max} = 144\pi$ cm/s. (Vì): Vận tốc cực đại $v_{max} = A\omega = 8 \times 6\pi = 48\pi$ cm/s, không phải $144\pi$ cm/s
\end{itemchoice}
}
\end{ex}

% ===== Câu 138 =====
% ID: L11C1B2-92-TLN
% Mức: VD
\begin{ex}
% ID: L11C1B2-92-TLN
% Mức: VD
Một vật dao động điều hòa theo phương trình $x = 6\cos\left(4\pi t - \dfrac{\pi}{3}\right)\,\text{cm}$, trong đó $t$ tính bằng giây. Tính vận tốc cực đại của vật theo đơn vị cm/s.
\shortans{$v_{max} = 24\pi \approx 75,4\,\text{cm/s}$}
\loigiai{
Từ phương trình $x = 6\cos\left(4\pi t - \dfrac{\pi}{3}\right)\,\text{cm}$, ta xác định được: biên độ $A = 6\,\text{cm}$, tần số góc $\omega = 4\pi\,\text{rad/s}$. Vận tốc cực đại của dao động điều hòa được tính theo công thức: $v_{max} = \omega A = 4\pi \times 6 = 24\pi \approx 75,4\,\text{cm/s}$.
}
\end{ex}

% ===== Câu 139 =====
% ID: L11C1B2-93-TLN
% Mức: VD
\begin{ex}
% ID: L11C1B2-93-TLN
% Mức: VD
Một vật dao động điều hòa theo phương trình $x = 5\cos\left(10\pi t + \dfrac{\pi}{4}\right)\,\text{cm}$, trong đó $t$ tính bằng giây. Tính gia tốc cực đại của vật theo đơn vị $\text{m/s}^2$.
\shortans{$a_{max} = 5\pi^2 \approx 49,3\,\text{m/s}^2$}
\loigiai{
Từ phương trình $x = 5\cos\left(10\pi t + \dfrac{\pi}{4}\right)\,\text{cm}$, ta xác định được: biên độ $A = 5\,\text{cm} = 0,05\,\text{m}$, tần số góc $\omega = 10\pi\,\text{rad/s}$. Gia tốc cực đại của dao động điều hòa được tính theo công thức: $a_{max} = \omega^2 A = (10\pi)^2 \times 0,05 = 100\pi^2 \times 0,05 = 5\pi^2 \approx 49,3\,\text{m/s}^2$.
}
\end{ex}

% ===== Câu 140 =====
% ID: L11C1B2-94-TLN
% Mức: VD
\begin{ex}
% ID: L11C1B2-94-TLN
% Mức: VD
Một vật dao động điều hòa theo phương trình $x = 8\cos\left(6\pi t + \dfrac{\pi}{6}\right)\,\text{cm}$, trong đó $t$ tính bằng giây. Tại thời điểm $t = \dfrac{1}{9}\,\text{s}$, vận tốc của vật có giá trị bằng bao nhiêu cm/s?
\shortans{$-24\pi\,\text{cm/s}$}
\loigiai{
Phương trình li độ: $x = 8\cos\left(6\pi t + \dfrac{\pi}{6}\right)\,\text{cm}$. Vận tốc là đạo hàm của li độ theo thời gian: $v = x' = -8 \cdot 6\pi \sin\left(6\pi t + \dfrac{\pi}{6}\right) = -48\pi \sin\left(6\pi t + \dfrac{\pi}{6}\right)\,\text{cm/s}$. Tại $t = \dfrac{1}{9}\,\text{s}$: $6\pi \cdot \dfrac{1}{9} + \dfrac{\pi}{6} = \dfrac{6\pi}{9} + \dfrac{\pi}{6} = \dfrac{2\pi}{3} + \dfrac{\pi}{6} = \dfrac{4\pi}{6} + \dfrac{\pi}{6} = \dfrac{5\pi}{6}$. Vậy $v = -48\pi \sin\left(\dfrac{5\pi}{6}\right) = -48\pi \cdot \dfrac{1}{2} = -24\pi\,\text{cm/s}$.
}
\end{ex}
